% Generated draft: requires mathematical and source review.
\documentclass[11pt]{article}
\usepackage[margin=25mm]{geometry}
\usepackage{fontspec}
\setmainfont{Latin Modern Roman}
\usepackage{amsmath,amssymb,mathtools,unicode-math}
\setmathfont{Latin Modern Math}
\usepackage{enumitem,needspace,xurl,hyperref}
\setlength{\parindent}{0pt}
\setlength{\parskip}{5pt}
\setlength{\emergencystretch}{3em}
\newcommand{\N}{\mathbb{N}}
\newcommand{\Z}{\mathbb{Z}}
\newcommand{\Q}{\mathbb{Q}}
\newcommand{\R}{\mathbb{R}}
\newcommand{\C}{\mathbb{C}}
\newcommand{\F}{\mathbb{F}}
\newcommand{\A}{\mathbb{A}}
\newcommand{\PP}{\mathbb{P}}
\newcommand{\E}{\mathbb{E}}
\newcommand{\eps}{\varepsilon}
\newcommand{\dd}{\,\mathrm{d}}
\begin{document}
\section*{UnsolvedMath ID: 20003270}
\textbf{Source problem code:} catalogue entry
\subsection{Definitions and mathematical statement}
The source is a boundary-damaged context paragraph, not the statement of Conjecture 4. According to the supplied aid, its opening continues ``Some partial results on Conjecture 4''. It attributes such partial results to [93, 142], F. Rodriguez-Hertz [162], and B. Kalinin and R. Spatzier. It then proposes broader classification directions:
\begin{itemize}
\item ``irreducible'' partially hyperbolic actions of $\Z^k$ and $\R^k$, $k\ge2$;
\item higher-rank actions of commuting expanding maps;
\item analogous actions of higher-rank connected semisimple Lie groups and their lattices under unspecified hyperbolicity assumptions, referring to [85, 131].
\end{itemize}
An action of $\Z^k$ is given by $k$ commuting invertible maps, whereas commuting expanding coverings generally define semigroup actions and need not be invertible. The paragraph does not define irreducibility, the precise partially hyperbolic or higher-rank hypotheses, the ambient manifolds, or the equivalence relation for classification. It therefore does not determine a single fully quantified classification conjecture.

The attached research aid reports this precise circle model. Write $\mathbb T=\R/\Z$. Let $f:\mathbb T\to\mathbb T$ be an orientation-preserving $C^1$ expanding covering of degree $m\ge2$, where expansion means that $|(f^j)'(x)|\ge C\lambda^j$ for all $x$ and $j\ge1$, for some $C>0$ and $\lambda>1$. If $g:\mathbb T\to\mathbb T$ is any continuous map of degree $n\in\Z$ with $fg=gf$, there are a homeomorphism $h:\mathbb T\to\mathbb T$ and $c\in\mathbb T$ such that
\[
 hfh^{-1}(x)=mx,\qquad hgh^{-1}(x)=nx+c,
 \qquad (m-1)c=0\quad\text{in }\mathbb T.
\]
Thus the residual translation phase lies in the finite group
\[
 \left\{\frac{j}{m-1}\pmod1:0\le j<m-1\right\}.
\]
The conclusion is a joint topological normal form; expansion of $g$ is not assumed, and smoothness of $h$ is not asserted.

The aid reports substantial progress by Spatzier and Yang for genuinely higher-rank smooth semigroup actions, without supplying their complete theorem or citation. It also cautions that general partially hyperbolic classification depends on hypotheses and that a blanket $C^1$ smooth-rigidity assertion is false. The circle model neither supplies those missing hypotheses nor classifies all the actions mentioned in the paragraph.
\subsection{Short English statement}
This fragment lists higher-rank classification directions rather than a complete conjecture. Its supplied concrete model jointly conjugates an expanding circle covering and any commuting continuous circle map to $x\mapsto mx$ and $x\mapsto nx+c$, with the exact finite constraint $(m-1)c=0$.
\subsection{Sources}
\begin{itemize}
\item American Institute of Mathematics, \emph{Emerging applications of measure rigidity}, context after Conjecture 4 and the transition to section 2.3, \url{https://aimath.org/WWN/measrigid/measrigid.pdf}. Supplied attributions: [93, 142], F. Rodriguez-Hertz [162], B. Kalinin, R. Spatzier, and [85, 131]; the bibliography entries are not included in this record.
\item The attached research summary attributes progress on genuinely higher-rank expanding semigroup actions to Spatzier and Yang but supplies no full bibliographic reference.
\item Local UnsolvedMath registry, permanent ID $20003270$, source code \texttt{catalogue entry}, canonical location \texttt{aim-other-notes.json}, \texttt{notes[76]}, dated 14 August 2026. Research attempt 1 is labelled context only, ``not a problem'' at run, proved rigor and low novelty confidence; it is a machine-generated, non-peer-reviewed aid requiring independent verification.
\end{itemize}
\end{document}
